\subsubsection{Training-time comparison}
\label{subsubsection:training_time_comparison}

We compare the fixed 192-dimensional NLDisco, CEBRA-Time, and LangevinFlow
configurations used in the Churchland reaching analysis on the same NVIDIA
H100 80-GB GPU. Each method is trained from scratch five times (seeds 0--4),
for 15 sequential fits in separate processes. Method order rotates between
seeds. All fits use the same 130,507 complete within-trial windows, each
containing ten 50-ms bins from 191 channels, with batch size 1,024.
NLDisco and LangevinFlow make ten exhaustive shuffled passes (1,280 updates).
CEBRA-Time samples temporal-positive reference pairs and negatives with
replacement: its ten epoch-equivalents comprise 1,275 updates and 1,305,600
reference draws, rather than exhaustive passes through the windows.
The fixed configurations are NLDisco with nested widths 96/192 and sparsity
budgets 8/16; CEBRA-Time with a 192-dimensional offset10 encoder, temperature
1 and lag 1; and LangevinFlow with state width 64, exposing 192 concatenated
position, velocity and recurrent-state coordinates, learning rate 0.003 and
a 500-epoch KL/weight-decay ramp. Preprocessing and objectives remain
method-specific, as detailed in the Churchland experiment.

Wall-clock timing is synchronized with CUDA immediately before and after
optimization. It includes minibatch preparation, device transfers during
training, loss checks and the trainer's local progress bookkeeping; it
excludes process startup, initial data loading/preprocessing, model and
optimizer construction, checkpoint writing, and latent export. Each fit
starts in a fresh process, without discarded training warm-up updates.
NLDisco's post-training threshold calibration is timed separately and
excluded from the plotted optimization time. Its median calibration time
is 3.61~s. External experiment tracking is disabled, and all
fits use four CPU threads and PyTorch 2.5.1 with CUDA 12.1. GPU process
occupancy is sampled every two seconds; no competing compute process is
observed on the benchmark GPU. Other GPUs on the host are not reserved.

\begin{figure}[htbp]
    \centering
    \includegraphics[width=0.8\linewidth]{figures/training_time_h100.pdf}
    \caption{\textbf{LVM training time comparison.}
    Models are fitted to the Churchland--Shenoy reaching recording of
    Nitschke from 12 August 2009 (1,906 valid trials; 191 recorded channels).
    One exhaustive epoch contains 130,507 overlapping, within-trial examples,
    each comprising ten consecutive 50-ms bins (500 ms of neural data;
    stride one bin). All three models expose a \textbf{192-dimensional latent
    representation}: NLDisco's largest dictionary (with a nested 96-dimensional
    level), CEBRA-Time's embedding, and LangevinFlow's concatenated position,
    velocity and recurrent states (64 dimensions each).
    Total/trainable parameter counts, excluding buffers, are
    NLDisco, 614,859/614,859; CEBRA-Time, 270,784/270,784; LangevinFlow, 449,199/449,199. Each dot is one fresh fit; five seeds (0--4) per method
    give 15 fits run sequentially on the same H100 80-GB GPU, with batch size
    1,024 and four CPU threads. NLDisco and LangevinFlow train for ten exhaustive
    epochs (1,280 updates). CEBRA-Time uses ten reference-draw epoch-equivalents
    with replacement sampling (1,275 updates; 1,305,600 reference draws in total),
    rather than exhaustive passes. CUDA-synchronized wall-clock time includes
    optimization and minibatch preparation, but excludes initial data loading
    and preprocessing, model/optimizer construction, checkpoint writing,
    latent export, and NLDisco threshold calibration (median 3.61~s).
    Boxes span the 25th--75th percentiles, horizontal lines and numeric labels
    mark medians, and whiskers reach the most extreme observations within
    1.5 interquartile ranges. All five observations per method are shown.
    Equal latent width does not imply equal model size, objectives, or
    convergence.}
    \label{figure:training_time_h100}
\end{figure}

\begin{center}
\begin{tabular}{lrr}
\toprule
Method & Median (s) & Interquartile interval (s) \\
\midrule
NLDisco & 43.42 & [43.36, 44.10] \\
CEBRA-Time & 8.18 & [6.51, 8.23] \\
LangevinFlow & 132.77 & [132.22, 132.89] \\
\bottomrule
\end{tabular}
\end{center}

These repeated timings compare the specified implementations at a fixed
data-exposure budget on common hardware. Equal exposed dimension does not
equalize parameter counts, training objectives, sampled observations, or
convergence; the measurements do not establish time to matched predictive
performance or discovery quality. The five fits characterize runtime and
initialization variability on one recording, not between-dataset variation.
